\section*{Chapter \ref{ch:b2:waves-on-strings} --- Waves on Strings and Rods: the d'Alembert Equation}

\begin{solution}{exo:b2:waves-on-strings:1}
$c = 2Lf = \qty{143}{m/s}$; $T = \mu c^2 = \qty{61}{N}$; $\lambda = 2L = \qty{1.3}{m}$ on the
string, $340/110 = \qty{3.1}{m}$ in air; fifth fret: $110 \times 2^{5/12} =
\qty{147}{Hz}$.
\end{solution}

\begin{solution}{exo:b2:waves-on-strings:2}
$\sqrt{E/\rho}$: steel \qty{5.1}{km/s}, aluminium \qty{5.1}{km/s}, lead
\qty{1.2}{km/s}. \qty{1}{km}: \qty{0.20}{s} in the rail, \qty{2.9}{s} in air ---
two sounds, the rail's first, \qty{2.7}{s} apart.
\end{solution}

\begin{solution}{exo:b2:waves-on-strings:3}
$400$, $600$, $800$ Hz with $1$, $2$, $3$ interior nodes; mode 3 nodes at
$L/3$ and $2L/3$. $T \times 2$: $200\sqrt2 = \qty{283}{Hz}$; $L/2$: \qty{400}{Hz}; $\mu
\times 3$: $200/\sqrt3 = \qty{115}{Hz}$.
\end{solution}

\begin{solution}{exo:b2:waves-on-strings:4}
$c = \sqrt{45/0.2} = \qty{15}{m/s}$; \qty{24}{m} in \qty{1.6}{s}; the pulse comes
back inverted ($r = -1$); with the free ring, upright ($r = +1$).
\end{solution}

\begin{solution}{exo:b2:waves-on-strings:5}
(a) $\mu = \rho\pi d^2/4 = \qty{6.1}{g/m}$; $c = 2Lf = \qty{334}{m/s}$; $T = \mu c^2 =
\qty{690}{N}$. (b) $\sigma = T/A = \rho c^2 = \qty{0.87}{GPa}$, margin $2.3$. (c) $T/4 =
\qty{170}{N}$, the same stress (it depends only on $\rho c^2$). (d) A long
string does not fit the case; a thick plain wire is too stiff (strongly
inharmonic) and hard to bend over the bridge; copper winding adds mass
without stiffness.
\end{solution}

\begin{solution}{exo:b2:waves-on-strings:6}
(a) $m_n = 4L^2f^2\mu/n^2g = 6.6/n^2$ kg: $6.6$, $1.65$, $0.73$, \qty{0.41}{kg}. (b)
$n^2 = 13.2$, $n = 3.6$: no resonance, between modes 3 and 4 (nearest: 4,
\qty{0.41}{kg}). (c) Practically a node: the driver's amplitude is tiny
compared with the loops. (d) Energy is fed in phase at each period and
accumulates; damping (air, internal friction, leakage through the
pulley) caps the amplitude.
\end{solution}

\begin{solution}{exo:b2:waves-on-strings:7}
$c = \qty{100}{m/s}$, $k = \qty{6.3}{rad/m}$, $\lambda = \qty{1.0}{m}$; $v_{\max} = \omega A =
\qty{1.3}{m/s}$, $a_{\max} = \omega^2A = \qty{790}{m/s^2}$; slope $kA = 0.013$; $Z =
\sqrt{T\mu} = \qty{1.0}{kg/s}$; $\langle\mathcal P\rangle = \tfrac12Z\omega^2A^2 = \qty{0.79}{W}$;
per wavelength $\mathcal P/f = \qty{7.9}{mJ}$.
\end{solution}

\begin{solution}{exo:b2:waves-on-strings:8}
(a) $c_1 = \qty{100}{m/s}$, $c_2 = \qty{50}{m/s}$; $Z_1 = \qty{0.20}{kg/s}$, $Z_2 = \qty{0.40}{kg/s}$.
(b) From the light string: $r = -1/3$, $t = 2/3$; from the heavy one: $r =
+1/3$, $t = 4/3$. (c) $R = 1/9$, $\mathcal T = 4 \times 0.08/0.36 = 8/9$ in both cases.
(d) Transmitted: \qty{0.67}{cm} high, \qty{5}{cm} wide (same duration, half
the speed); reflected: \qty{-0.33}{cm}, \qty{10}{cm}.
\end{solution}

\begin{solution}{exo:b2:waves-on-strings:9}
(a) $E = F\ell/A\Delta\ell = 80 \times 2/(7.85 \times 10^{-7} \times 10^{-3}) = \qty{2.0e11}{Pa}$.
(b) $c = \qty{5.1}{km/s}$. (c) Longitudinal: $c/2\ell = \qty{1.3}{kHz}$; transverse:
$\sqrt{T/\mu} = \qty{114}{m/s}$, $f = \qty{29}{Hz}$ --- the tension stresses the
wire to \qty{100}{MPa}, $5 \times 10^{-4}$ of $E$, so the transverse speed is
$\sqrt{5 \times 10^{-4}}$ of the longitudinal one. (d) $k = Ea = \qty{50}{N/m}$; $m
= \rho a^3 = \qty{1.2e-25}{kg}$ (about $70$ atomic mass units).
\end{solution}

\begin{solution}{exo:b2:waves-on-strings:10}
(a) $b_n = \frac2L\int_0^Ly(x, 0)\sin(n\pi x/L)\dd x$ with the triangle $y =
2hx/L$ on $[0, L/2]$ (and symmetric): integrate by parts, $b_n = 8h\sin(n\pi/2)
/n^2\pi^2$. (b) Even $n$ vanish: the midpoint, a node of every even mode,
is the point pulled. (c) $E_n = \tfrac14\mu L\omega_n^2b_n^2 \propto n^2/n^4 = 1/n^2$
(odd $n$): the fundamental has $1/\sum_{\text{odd}}n^{-2} = 8/\pi^2 = 81\%$. (d)
$b_n \propto \sin(n\pi/5)/n^2$: $n = 5, 10, \dots$ vanish. Near the bridge the
factor $\sin(n\pi x_0/L)$ is small for the low modes and keeps growing
for the higher ones: relatively more high harmonics, a brighter sound.
\end{solution}

\begin{solution}{exo:b2:waves-on-strings:11}
(a) $-m\omega^2 = k(2\cos ka - 2) = -4k\sin^2(ka/2)$. (b) $v_\varphi = 2\omega_0\sin(ka/2)
/k$, $v_g = \omega_0a\cos(ka/2)$; both $\to \omega_0a = c$ as $ka \to 0$; at $ka = \pi$:
$v_\varphi = 2c/\pi$, $v_g = 0$. (c) Neighbours move in antiphase: a \omterm{prop:b2:waves-on-strings:modes}{standing
wave} that transports nothing; a $k$ beyond $\pi/a$ gives the same set of
displacements as some $k' = k - 2\pi/a$ --- no new wave. (d) $\omega_0 = c/a =
\qty{2.0e13}{rad/s}$, $f_{\max} = \omega_0/\pi = \qty{6.5}{THz}$; \qty{1}{MHz}: $\lambda =
\qty{5.1}{mm} = 2 \times 10^7a$, $ka \sim 3 \times 10^{-7}$: perfectly non-dispersive.
\end{solution}

\begin{solution}{exo:b2:waves-on-strings:12}
(a) $T(x) = \mu gx$, $c = \sqrt{gx}$. (b) $t = \int_0^L\dd x/\sqrt{gx} = 2\sqrt{L/g} =
\qty{2.0}{s}$; the upper part of a pulse moves faster than the lower: it
stretches. (c) $c = \sqrt{T/\mu}$ grows toward the tip. (d) Energy $\approx
\mathcal P\tau = ZV^2\tau = \sqrt{T\mu}V^2\tau$ constant: $V \propto \mu^{-1/4}$; a
factor $10^4$ in $\mu$ gives $\times10$: \qty{100}{m/s} --- a real whip gains
the rest from the unrolling loop, which concentrates the energy
further, and the tip passes the speed of sound.
\end{solution}

\begin{solution}{pb:b2:waves-on-strings:1}
\textbf{1.} $\mu = \rho\pi d^2/4 = \qty{6.1}{g/m}$; $c = 2Lf = \qty{334}{m/s}$; $T = \mu c^2
= \qty{685}{N}$.

\textbf{2.} $\sigma = \rho c^2 = \qty{0.87}{GPa}$: margin $2.3$. $T \propto f^2$: a
semitone sharp adds $12\%$, an octave too high quadruples it,
\qty{3.5}{GPa} --- it snaps.

\textbf{3.} $230 \times 685 \approx \qty{160}{kN}$, sixteen tonnes.

\textbf{4.} $c = 2Lf = \qty{105}{m/s}$; at \qty{685}{N}: $\mu = T/c^2 = \qty{63}{g/m}$,
i.e.\ a steel rod of \qty{3.2}{mm} --- far too stiff to vibrate
harmonically or to bend over the bridge; a thin core wound with copper
gives the mass without the stiffness.

\textbf{5.} $440$, $880$, $1320$, $1760$, $2200$, $2640$, $3080$, \qty{3520}{Hz}: $2$,
$4$, $8$ are octaves; $3$, $6$ a fifth above an octave; $5$ a just major
third above two octaves; $7$ ($3080/1760 = 1.75$) is a flat minor
seventh --- dissonant.

\textbf{6.} $E = \tfrac14\mu L\omega^2A^2 = 0.25 \times 6.1 \times 10^{-3} \times 0.38 \times (2\pi
\times 440)^2 \times (5 \times 10^{-4})^2 = \qty{1.1}{mJ}$.

\textbf{7.} $A\sin kx\cos\omega t = \tfrac12A[\sin(kx - \omega t) + \sin(kx + \omega t)]$: two
waves of amplitude \qty{0.25}{mm}; $Z = \mu c = \qty{2.0}{kg/s}$; each carries
$\tfrac12Z\omega^2(A/2)^2 = \qty{0.49}{W}$.

\textbf{8.} $\dot y(x, 0) = \sum_nB_n\omega_n\sin(n\pi x/L)$; multiply by
$\sin(m\pi x/L)$ and integrate, using $\int_0^L\sin\sin = \tfrac L2\delta_{mn}$.

\textbf{9.} The integrand is $v_0\sin(n\pi x/L)$ over a width $w$ around $x_0$,
where the sine is nearly constant.

\textbf{10.} $\sin(n\pi/8) = 0$ for $n = 8, 16, \dots$; and since $B_n/B_1 =
\sin(n\pi/8)/(n\sin(\pi/8))$, the seventh harmonic is reduced to $1/7$ of the
first while the consonant $2$ to $6$ keep $0.9$ to $0.3$: the dissonant
seventh is tamed.

\textbf{11.} Two velocity pulses leave the struck point in opposite
directions, reflect inverted at the ends, cross and recombine; the
initial state recurs after $2L/c = 1/f_1 = \qty{2.3}{ms}$.

\textbf{12.} $B = \pi^3 \times 2 \times 10^{11} \times 10^{-12}/(64 \times 685 \times 0.144) =
9.8 \times 10^{-4}$; $\sqrt{1 + 64B} = 1.031$: $+3.1\%$, $1200\log_2(1.031) = 53$ cents.

\textbf{13.} The harmonics of a low note are sharp of the exact
multiples; to make them coincide with the fundamentals of the higher
notes (and avoid beats), the high notes are tuned slightly sharp.

\textbf{14.} $Z_s = \mu c = \qty{2.0}{kg/s}$; $r = (2 - 2000)/2002 = -0.998$.

\textbf{15.} $1 - r^2 \approx 4Z_s/Z_b = 4 \times 10^{-3}$.

\textbf{16.} $c/2L = 440$ reflections per second: decay rate $440 \times 4 \times
10^{-3} = \qty{1.8}{s^{-1}}$, $\tau = \qty{0.56}{s}$; $\ln10^6/1.8 = \qty{7.7}{s}$ for
\qty{60}{dB}.

\textbf{17.} $1.1\,\text{mJ} \times 1.8\,\text{s}^{-1} = \qty{2}{mW}$ --- two hundred
times a quiet conversation; it is what you hear.

\textbf{18.} A \qty{1}{mm} wire pushes almost no air: the air slips round
it and its radiated power is negligible. The soundboard takes the
string's energy through the bridge and moves air over a square metre:
an impedance match by area.

\textbf{19.} \qty{440}{Hz} is the second harmonic of A$_3$: the bridge drives
that string at one of its modes --- sympathetic resonance.

\textbf{20.} Raise $Z_b$: less energy leaks per reflection, the note lasts
longer --- and is quieter.

\textbf{21.} $c = \qty{5.1}{km/s}$: \qty{1.0}{s} in the rail, \qty{15}{s} in air.

\textbf{22.} $50/5064 = \qty{10}{ms}$; the stress must vanish at a free end:
a compression returns as a rarefaction.

\textbf{23.} $\lambda = \qty{2.5}{mm}$; $Z_{\text{steel}} = \rho c = \qty{4e7}{kg.m^{-2}.s^{-1}}$
against $400$: $r \approx -1$, total reflection: any open crack, however
thin, returns an echo.

\textbf{24.} $\Delta t = d(1/3.5 - 1/6) = d \times \qty{0.119}{s/km}$: $d = \qty{170}{km}$.

\textbf{25.} String \qty{334}{m/s} ($T$, $\mu$); rail \qty{5.1}{km/s} ($E$, $\rho$);
rock P \qty{6}{km/s} (compression modulus, $\rho$); rock S \qty{3.5}{km/s}
(shear modulus, $\rho$); air \qty{340}{m/s} ($\gamma P$, $\rho$ --- next chapter).
\end{solution}
